% DOCUMENT FORMAT ======================= -*- mode: LaTeX; coding: utf-8 -*- ===

\documentclass[diploma]{softlab-thesis}


% PACKAGE SETTINGS =============================================================

\usepackage[cm-default]{fontspec}
\usepackage{amsmath}
\usepackage{amsfonts}
\usepackage{multirow}
\usepackage{array}
\usepackage{mdwlist}
\usepackage{subfig}
\usepackage{floatrow}
%\usepackage{float}
\usepackage{verbatim}
\usepackage{color}
\usepackage{graphicx}
\usepackage{xunicode}
\usepackage{xltxtra}
\usepackage{url}
%\usepackage{dsfont}
%\usepackage{microtype}
\usepackage[htt]{hyphenat}
\usepackage{multicol}
\usepackage{wrapfig}

% FONT SETTINGS ===============================================================

%\setromanfont[Mapping=tex-text]{CMU Serif}
%%\setromanfont[Mapping=tex-text]{CMU Sans Serif} % temporary change until printing
%%\setsansfont[Mapping=tex-text]{CMU Sans Serif}
%%\setmonofont[Mapping=tex-text]{CMU Typewriter Text}
%\setmainfont[Mapping=tex-text]{CMU Serif}
%%\setmainfont[Mapping=tex-text]{CMU Sans Serif}  % temporary change until printing

%\setromanfont[Mapping=tex-text,ExternalLocation=fonts/]{cmunrm.otf}
%\setsansfont[Mapping=tex-text,ExternalLocation=fonts/]{cmunss.otf}
%\setmonofont[Mapping=tex-text,ExternalLocation=fonts/]{cmuntt.otf}
%\setmainfont[Mapping=tex-text, ExternalLocation=fonts/]{cmunss.otf}

% CUSTOM COMMANDS ==============================================================

% Layout macros
\newcommand\spa[1]{\; #1 \;}

% Font macros
\newcommand\resfont[1]{\ensuremath{\mathtt{#1}}}

% Mathematical macros
\newcommand\setmap[3]{#1\{#2 \mapsto #3\}}
\newcommand\getmap[3]{(#2 \mapsto #3) \in #1}
\newcommand\tuple[2]{\ensuremath{\langle#1, #2\rangle}}
\newcommand\mfrac[2]{\ensuremath{\dfrac{#1}{#2}}}
\newcommand\nequiv[2]{\ensuremath{#1 \not\equiv #2}}

% Core Ruby Operational Semantics letter bindings
\newcommand\mem{\mu}

% Core Ruby Operational Semantics low level macros
\newcommand\state[2]{(#1, #2)}
\newcommand\transition[1]{\ensuremath{\overset{#1, c*}{\rightarrow}}}
\newcommand\range[2]{#1, ..., #2}
\newcommand\midrange[5]{\range{#1}{#2}, #3, \range{#4}{#5}}

% Core Ruby Operational Semantics high level macros
\newcommand\operation[5]{\ensuremath{\state{#1}{#2} \transition{#3} \state{#4}{#5}}}
\newcommand\propagation[2]{\operation{#1}{\mem}{#2}{#1'}{\mem'}}

% Core Ruby specific Operational Semantics macros
\newcommand\semicolon[2]{#1; \; #2}
\newcommand\assign[2]{#1 = #2}
\newcommand\mcall[3]{#1.\texttt{#2}(#3)}
\newcommand\ifte[3]{\resfont{if} \; #1 \; \resfont{then} \; #2 \; \resfont{else} \; #3}
\newcommand\newclass[2]{#1.\resfont{new}(#2)}
\newcommand\methoddef[3]{\resfont{def} \; #1(#2) = #3}
\newcommand\classdef[2]{\resfont{class} \; #1 = #2}
\newcommand\with[3]{with \; \tuple{#1}{#2} \; do \; #3}

% Success Typing macros
\newcommand\ssub{\sqsubseteq_S}

% Core Ruby Success Typing letter bindings
\newcommand\classlist{\Delta}
\newcommand\envir{\Gamma}
\newcommand\fields{\Phi}
\newcommand\currclass{l}

% Core Ruby Success Typing inferencing macros
\newcommand\stinfer[5]{\classlist; \; #1; \; \fields \; \underset{\currclass}{\vdash} \; #2: #3 \; \& \; #4; \; #5}


%%%%%%%%%%%%%%%%%%%%%%%%%%% HASKELL STUFF %%%%%%%%%%%%%%%%%%%%%%%%%%%

%\newcommand\typerep[1]{\ensuremath{#1}}
\newcommand\typerep[1]{\lstinline[basicstyle=\normalsize\ttfamily,keywords={}]|#1|}
\newcommand\typefootrep[1]{\textbf{\lstinline[basicstyle=\footnotesize\ttfamily,keywords={}]|#1|}}
% \newcommand\ttyperep[1]{\typerep{#1}}
% \newcommand\mtyperep[1]{\mbox{\typerep{#1}}}

% Arrow types
\newcommand\typeto[2]{\typerep{#1} \typerep{->} \typerep{#2}}
\newcommand\typetotwo[3]{\ensuremath{\typerep{#1} \typerep{->}
                                     \typerep{#2} \typerep{->}
                                     \typerep{#3}}}

\newcommand\tyconapone[2]{\ensuremath{\mbox{\typerep{#1}} \:\: \mbox{#2}}}
\newcommand\tyconaponeC[2]{\ensuremath{\mbox{\typerep{#1}} \:\: \mbox{\typerep{#2}}}}
% \newcommand\tyconapone[2]{\typerep{#1} $\:$ \typerep{#2}}

\newcommand\tyconaptwo[3]{\ensuremath{\mbox{\typerep{#1}} \:\: \mbox{#2} \:\: \mbox{#3}}}


% FLOAT SETUP ==================================================================

% SPELLING =====================================================================

% CODE HIGHLIGHTING ============================================================

\usepackage{listings}

\definecolor{gray}{rgb}{0.5,0.5,0.5}
\definecolor{darkgreen}{rgb}{0.0,0.5,0.0}

\lstset{
	language=Haskell,
	basicstyle=\footnotesize\ttfamily,
	numbers=left,
	numberstyle=\tiny\color{gray},
	stepnumber=1,
	numbersep=5pt,
	backgroundcolor=\color{white},
	showspaces=false,
	showstringspaces=false,
	showtabs=false,
	frame=single,
	rulecolor=\color{black},
	tabsize=2,
	captionpos=b,
	breaklines=true,
	breakatwhitespace=false,
	title=\lstname,
        xleftmargin=17pt,
        xrightmargin=5pt,
        aboveskip=\baselineskip
}

\newcommand\haskellcode[2]{\lstinputlisting[language=Haskell,caption={#1},label=#1]{src/#2}}
\newcommand\rubycode[2]{\lstinputlisting[language=Ruby,caption={#1},label=#1]{src/#2}}
\newcommand\erlangcode[1]{\lstinputlisting[language=erlang,captionpos=]{src/#1}}
\newcommand\plaintext[1]{\lstinputlisting[language=,captionpos=,stepnumber=0]{src/#1}}

% CHANGE MATH FONT ============================================================

% HYPERREF MUST BE LAST =======================================================

\usepackage[xetex,colorlinks=true,linkcolor=blue,citecolor=darkgreen]{hyperref}

% DOCUMENT INFORMATION =========================================================

\title{Thesis subject} % ===============> FIXME
\author{Όνομα Φοιτητή}
\authoren{Student name}
\datedefense{9}{9}{9999}      % ===============> FIXME
\TRnumber{CSD-SW-TR-*-*} % ok ==================> FIXME
\supervisor{Υπέυθυνος Διπλωματικής}
\supervisorpos{Τίτλος Υπευθύνου}
\committeeone{Πρώτο μέλος επιτροπής}
\committeeonepos{Τίτλος μέλους}
\committeetwo{Δεύτερο μέλος επιτροπής}
\committeetwopos{Τίτλος μέλους}
\committeethree{Τρίτο μέλος επιτροπής}
\committeethreepos{Τίτλος μέλους}

\hypersetup{
	pdftitle={Thesis subject}, % ===> FIXME
	pdfauthor={Όνομα Φοιτητή},
	pdfsubject={},
	pdfkeywords={}
}


% MAIN DOCUMENT ================================================================

\begin{document}

\frontmatter
\maketitle

\def\templen{\parindent}
\setlength{\parindent}{0pt}
\setlength{\parskip}{1.5ex plus 0.5ex minus 0.2ex}
% DOCUMENT FORMAT ======================= -*- mode: LaTeX; coding: utf-8 -*- ===
\begin{abstractgr}
	Περίληψη της διπλωματικής.
\begin{keywordsgr}
  	Λέξη-κλειδί 1, λέξη-κλειδί 2, λέξη-κλειδί 3
\end{keywordsgr}
\end{abstractgr}

\begin{abstracten}
	Abstract of diploma thesis.
\begin{keywordsen}
  	Key-word 1, Key-word 2, Key-word 3
\end{keywordsen}
\end{abstracten}

\begin{acknowledgementsgr}
	Ευχαριστίες.
\end{acknowledgementsgr}





\selectlanguage{english}   % Everything from here on is in english
\setlanguage{english}
\nogreekalph

\setlength{\parindent}{\templen}
\setlength{\parskip}{0pt}
\tableofcontents
\listoffigures
%\listoftables
\renewcommand{\lstlistlistingname}{List of Listings}
%\lstlistoflistings % changed the title above

\mainmatter
% moved these two commands here so that they don't influence the toc
\setlength{\parindent}{0pt}
\setlength{\parskip}{1.5ex plus 0.5ex minus 0.2ex}

\renewcommand\floatpagefraction{.7}

\chapter{Chapter 1}\label{ch:ch1}

\section{Section 1}

First section.

\section{Section 2}

Second section.
It also includes a minor citation\cite{misc_citation}.

\section{Outline}

You can choose between two types of outlines. This:

Chapter~\ref{ch:ch2} does this and this. Chapter~\ref{ch:ch3} does that and
that. Chapter~\ref{ch:ch4} does these and these. And finally,
chapter~\ref{ch:ch5} does those and those.

or this:

\begin{description}
\item[Chapter]~\ref{ch:ch2}
	We do this and this.
\item[Chapter]~\ref{ch:ch3}
	We do that and that.
\item[Chapter]~\ref{ch:ch4}
	We do these and these.
\item[Chapter]~\ref{ch:ch5}
	We do those and those.
\end{description}

\chapter{Chapter 2}\label{ch:ch2}

\section{Section 1}

This section has an important citation\cite{inproceeding}

\subsection{Subsection 1}

This subsection has code in Haskell:
\haskellcode{Sample code}{samplecode.hs}

It also has a list:
\begin{description}
\item[Item 1]
	First item
\item[Item 2]
	Second item and a footnote\footnote[1]
	{Footnote description.}.
\item[Item 3]
	Third item and text in \textit{italics}.
\end{description}

And an enumerated list:

\begin{enumerate}
\item	First item.
\item	Second item and text in \textbf{bold}
\end{enumerate}

\section{Section 2}

\subsection{Subsection 1}
This subsection has a link to the block of code~\ref{Sample code} in Section 1.

\subsection{Subsection 2}

% FIXME: Must add content here
This subsection has a FIXME comment, visible only to the author.

\chapter{Chapter 3}\label{ch:ch3}

\section{Section 1}

This how we add a url: \url{http://www.example.org}

\section{Section 2}

And this is how we point to Figure~\ref{fig:image1}.

\begin{figure}[hb]
	\centering
	\includegraphics[]{diagrams/9.jpg}
	\caption{This is an image}
	\label{fig:image1}
\end{figure}

\subsection{Sub-section 3}

Another way to create a list:
\begin{itemize}
	\item Item 1
	\item Item 2
	\item Item 3
	\item Item 4
\end{itemize}


\chapter{Chapter 4}\label{ch:ch4}

\section{Section 1}

We can also use a special fonts to differentiate between text and $math()$.


\chapter{Chapter 5}\label{ch:ch5}

\section{Section 1}

This is an equation:
\begin{equation}
	true \in 9
\label{rule:9}
\end{equation}

which can once again be referenced like so \ref{rule:9}.


\backmatter
\cleardoublepage % start at the next odd page
\phantomsection % correct hyperlinking
\bibliography{references}
\bibliographystyle{softlab-thesis}
% \include{glossary}
% \chapter{Appendix}
% \printindex

\end{document}
